% ============================================================
% Autor: Eloi Egea Rada
% Tema: Complete specification of system, technological, and quality requirements
% ============================================================

\chapter{Functional and Non-Functional Requirements}
\label{ch:requirements}

This chapter specifies the requirements for the EVE-NG-based~\cite{eve_ng} cybersecurity
teaching package: functional requirements (RF) define what the system must do,
technological requirements (TR) specify the technical environment, and non-functional
requirements (NFR) address quality attributes and compliance.

\section{Functional Requirements}

\subsection{RF1: Design and Implement EVE-NG Topologies}

\textbf{Description}: Create four separate, functional, and interconnected network topologies
in EVE-NG format (.unl files) that simulate real-world cybersecurity scenarios.

\begin{itemize}
\item \textbf{RF1.1}: Four topologies, each targeting a specific pentesting phase:
  \begin{itemize}
  \item Lab 1: Reconnaissance (passive intelligence gathering)
  \item Lab 2: Web application security (active vulnerability discovery)
  \item Lab 3: Incident response and network forensics (lateral movement and log analysis)
  \item Lab 4: Privilege escalation (post-exploitation scenarios)
  \end{itemize}
\end{itemize}

\textbf{Success Criteria}: See technical KPIs in Chapter~3 (Table~\ref{tab:kpi_all}).

\subsection{RF2: Implement Vulnerabilities and Attack Scenarios}

\textbf{Description}: Implement realistic vulnerabilities aligned with PTES\cite{ptes},
OWASP\cite{owasp_testing_guide}, NIST\cite{nist_framework}, and CEH\cite{ceh_framework}.

\begin{itemize}
\item \textbf{RF2.1}: Lab 1 --- Reconnaissance: passive tools (Shodan, DNS enumeration),
      active scanning (Nmap\cite{nmap}, banner grabbing), service discovery and reporting.

\item \textbf{RF2.2}: Lab 2 --- Web Vulnerabilities (OWASP Top 10): SQL injection,
      command injection, authentication bypass, session flaws, XSS, CSRF,
      broken access control, and security misconfigurations.

\item \textbf{RF2.3}: Lab 3 --- Incident Response: lateral movement analysis,
      log forensics (auth.log, bash\_history), privilege escalation via sudo
      misconfiguration, network traffic analysis with Wireshark\cite{wireshark}.

\item \textbf{RF2.4}: Lab 4 --- Privilege Escalation (CEH): local privilege escalation,
      SUID/SGID binaries, sudo misconfigurations, writable cron jobs, capabilities,
      and post-exploitation persistence.
\end{itemize}

\textbf{Success Criteria}: Minimum 3--4 exploitable scenarios per lab, all resolvable
with documented remediation procedures.

\subsection{RF3: Create Structured Assessment Materials}

\textbf{Description}: Develop assessment rubrics and evaluation materials aligned with
the GEISI curriculum~\cite{geisi_curriculum}.

\begin{itemize}
\item \textbf{RF3.1}: 3--4 specific learning outcomes per lab, mapped to BLOOM's taxonomy
      (knowledge, application, analysis) and GEISI degree competencies.

\item \textbf{RF3.2}: Detailed assessment rubric per lab (100~pt scale with explicit
      criteria), included in the instructor resolution guide.
\end{itemize}

\textbf{Success Criteria}: 4 student lab guides and 4 instructor guides
(resolució PDF), each including a rubric.

\section{Technological Requirements}

\subsection{TR1: Virtualisation Infrastructure}

\begin{itemize}
\item \textbf{TR1.1}: EVE-NG 5.0+ (Community and Professional editions), deployed as a
      VM on Proxmox VE; web interface via HTTPS; RESTful API for automation.

\item \textbf{TR1.3}: Host specifications: 2$\times$ Intel Xeon E5-2680v4 (56 logical cores),
      64~GB RAM, 2~TB NVMe storage, Gigabit Ethernet.
\end{itemize}

\subsection{TR2: Security Tools and Penetration Testing Framework}

\begin{itemize}
\item \textbf{TR2.1}: Attacking platform: Parrot OS Security 6.x~\cite{parrot_os} (primary).
      Selected over Kali Linux for lower memory/disk footprint in concurrent EVE-NG
      deployments, cleaner QEMU integration, and a more approachable environment for
      students new to security distributions.

\item \textbf{TR2.2}: Core tools on Parrot OS: Nmap, Metasploit 6.x\cite{metasploit},
      Burp Suite\cite{burp_suite}, OWASP ZAP\cite{owasp_zap}, Wireshark\cite{wireshark},
      tcpdump, Python 3.10+, Bash.

\item \textbf{TR2.3}: Target images: Ubuntu Server 24.04~LTS, Ubuntu Desktop 24.04~LTS,
      pfSense~CE, VyOS, and the custom SYNAPSE web application (Lab~2).
\end{itemize}

\subsection{TR3: Scripting and Automation}

\begin{itemize}
\item \textbf{TR3.1}: Bash and Python 3.10+ for deployment, reset, and validation scripts.

\item \textbf{TR3.2}: At least one utility script per lab; ISO uploader
      (\texttt{Iso\_uploader.py}), bridge configuration, and connectivity check scripts;
      all documented with usage instructions in the repository.
\end{itemize}

\section{Non-Functional Requirements}

\begin{table}[H]
  \centering
  \small
  \setlength{\tabcolsep}{4pt}
  \renewcommand{\arraystretch}{1.15}
  \begin{tabularx}{\textwidth}{l l X}
    \toprule
    \textbf{NFR} & \textbf{Attribute} & \textbf{Key Requirements} \\
    \midrule
    NFR1 & Performance    & Deploy $\leq$2 min; VM boot $\leq$90 s (95\%); reset $\leq$30 s;
                            VM image $\leq$5 GB; $\leq$6 GB RAM per lab instance \\
    NFR2 & Compliance     & NIST CSF functions (Identify/Protect/Detect/Respond/Recover);
                            CIS Controls v8; OWASP Top~10 full coverage; PTES + CEH alignment \\
    NFR3 & Reliability    & MTBF $\geq$100 h; MTTR $\leq$5 min; 100\% scenario reproducibility;
                            zero critical bugs at release \\
    NFR4 & Security       & Full network isolation from production; no outbound internet from labs;
                            RBAC for lab management; mandatory ethics agreement \\
    NFR5 & Usability      & Faculty deploy time 5--10 min; student prerequisites minimal;
                            WCAG 2.1 Level~AA for web interfaces \\
    NFR6 & Sustainability & All components open-source or freely available; no vendor lock-in;
                            explicit mapping to GEISI learning outcomes; horizontally scalable \\
    \bottomrule
  \end{tabularx}
  \caption{Non-functional requirements summary}
  \label{tab:nfr_summary}
\end{table}

\section{Requirements Traceability Matrix}

\begin{table}[H]
  \centering
  \small
  \setlength{\tabcolsep}{4pt}
  \renewcommand{\arraystretch}{1.1}
  \begin{tabularx}{\textwidth}{l l X l}
    \toprule
    \textbf{ID} & \textbf{Type} & \textbf{Requirement} & \textbf{Scope} \\
    \midrule
    RF1.1  & Functional  & Design 4 EVE-NG topologies              & All labs       \\
    RF2.1  & Functional  & Lab 1 reconnaissance scenarios          & Lab 1          \\
    RF2.2  & Functional  & Lab 2 OWASP vulnerabilities             & Lab 2          \\
    RF2.3  & Functional  & Lab 3 incident response / forensics     & Lab 3          \\
    RF2.4  & Functional  & Lab 4 privilege escalation              & Lab 4          \\
    RF3.1  & Functional  & Learning outcomes defined               & All labs       \\
    RF3.2  & Functional  & Assessment rubrics                      & All labs       \\
    TR1.1  & Technical   & EVE-NG 5.0+                             & Platform       \\
    TR1.3  & Technical   & Host 64~GB RAM, 2~TB storage            & Infrastructure \\
    TR2.1  & Technical   & Parrot OS Security as attacking platform & Tools         \\
    NFR1   & Non-Func.   & Deploy $\leq$2 min, reset $\leq$30 s   & Performance    \\
    NFR2   & Non-Func.   & NIST / CIS / OWASP / PTES alignment    & Compliance     \\
    NFR4   & Non-Func.   & Network isolation + ethics              & Security       \\
    NFR6   & Non-Func.   & Open-source, GEISI aligned              & Sustainability \\
    \bottomrule
  \end{tabularx}
  \caption[Requirements traceability]{Requirements traceability matrix}
  \label{tab:requirements_traceability}
\end{table}

Detailed sub-requirements for RF1.2--RF1.4, RF2.5, RF3.3--RF3.4, TR1.2, TR1.4, TR2.4, TR3.3, and the full NFR subsections are provided in the \href{https://TheEgea.github.io/TFG/annexos/app-h/}{project online documentation}.
